% Authority: German Wikiversity pageid 151168, revid 1095913, 2026-06-15T07:25:24Z.
% Exact UTF-8 wikitext witness SHA-256: c771b7c1ed173ddc05c8c4219c3c34e7e8099fe42079740b292d991b2ce3fb25.
% Expanded German TeX source SHA-256: 7e3437274bf4f79b6a3fa719876b09fdbcee4e10523a215e9e6f4ecb566cdb85.

Nilai-nilainya adalah

\mavergleichskettedisp
{\vergleichskette
{ \gamma(0)
}
{ =} { \begin{pmatrix}  1  \\ 0   \end{pmatrix}
}
{ } {
}
{ } {
}
{ } {
}
}
{}{}{,}

\mavergleichskettedisp
{\vergleichskette
{ \gamma'(0)
}
{ =} { \begin{pmatrix}  0  \\ \omega   \end{pmatrix}
}
{ } {
}
{ } {
}
{ } {
}
}
{}{}{}
dan

\mavergleichskettedisp
{\vergleichskette
{ \gamma^{\prime \prime} (0)
}
{ =} { \begin{pmatrix}  - \omega^2  \\ 0   \end{pmatrix}
}
{ } {
}
{ } {
}
{ } {
}
}
{}{}{.}
Suatu gerak melingkar yang memiliki nilai serta turunan pertama dan kedua yang sama dengan gerak yang diberikan pada

\mavergleichskette
{\vergleichskette
{ t
}
{ = }{ 0
}
{  }{
}
{    }{
}
{   }{
}
}
{}{}{}
harus, khususnya, mempunyai garis singgung yang sama di titik ini. Karena itu, pusat lingkaran harus terletak pada sumbu $x$. Selain itu, pusat tersebut harus berada di sebelah kiri
\mathl{\begin{pmatrix}  1  \\ 0   \end{pmatrix}}{},
karena percepatan mengarah ke kiri. Oleh sebab itu, kita dapat menuliskan gerak melingkar beraturan yang dicari, dengan pusat
\mathl{\begin{pmatrix}  x  \\ 0   \end{pmatrix}}{}
\zusatzklammer {
\mavergleichskettek
{\vergleichskettek
{ x
}
{ < }{  1
}
{  }{
}
{    }{
}
{   }{
}
}
{}{}{}} {} {}
dan jari-jari
\mathl{1-x}{}, sebagai

\mavergleichskettedisp
{\vergleichskette
{  \delta (t)
}
{ =} {  \begin{pmatrix}  x  \\ 0   \end{pmatrix} + (1-x) \begin{pmatrix}  \cos \left(  u t \right)  \\ \sin  \left(  u t \right)     \end{pmatrix}
}
{ } {
}
{ } {
}
{ } {
}
}
{}{}{.}
Kita memperoleh

\mavergleichskettedisp
{\vergleichskette
{ \delta(0)
}
{ =} { \gamma(0)
}
{ } {
}
{ } {
}
{ } {
}
}
{}{}{.}
Selanjutnya,

\mavergleichskettedisp
{\vergleichskette
{ \delta'( 0)
}
{ =} {  (1-x)  \begin{pmatrix}  0  \\ u     \end{pmatrix}
}
{ } {
}
{ } {
}
{ } {
}
}
{}{}{}
dan

\mavergleichskettedisp
{\vergleichskette
{ \delta^{\prime \prime}( 0)
}
{ =} {  (1-x)  \begin{pmatrix}  -u^2  \\ 0     \end{pmatrix}
}
{ } {
}
{ } {
}
{ } {
}
}
{}{}{.}
Jadi, agar syarat-syarat yang diminta terpenuhi, harus berlaku

\mavergleichskettedisp
{\vergleichskette
{ \omega
}
{ =} { (1-x) u
}
{ } {
}
{ } {
}
{ } {
}
}
{}{}{}
dan

\mavergleichskettedisp
{\vergleichskette
{ \omega^2
}
{ =} { (1-x) u^2
}
{ } {
}
{ } {
}
{ } {
}
}
{}{}{.}
Dari sini diperoleh

\mavergleichskettedisp
{\vergleichskette
{ \omega^2
}
{ =} {  \omega u
}
{ } {
}
{ } {
}
{ } {
}
}
{}{}{}
dan dengan demikian

\mavergleichskettedisp
{\vergleichskette
{ u
}
{ =} { \omega
}
{ } {
}
{ } {
}
{ } {
}
}
{}{}{}
serta

\mavergleichskettedisp
{\vergleichskette
{ x
}
{ =} { 0
}
{ } {
}
{ } {
}
{ } {
}
}
{}{}{.}
